% !TeX root = ../main.tex
\section{Introducción}
\label{sec:introduccion}

El auge de la Inteligencia Artificial (IA) y del Aprendizaje Automático (\textit{Machine Learning}, ML) ha transformado sectores críticos como la salud, las finanzas y la conducción autónoma. No obstante, el entrenamiento de modelos de gran capacidad requiere volúmenes masivos de datos, frecuentemente sensibles o confidenciales (geolocalización, hábitos de consumo o registros biométricos). En los paradigmas tradicionales centralizados, la recopilación de información en un único repositorio introduce graves riesgos de privacidad, problemas de cumplimiento normativo (como el RGPD) y cuellos de botella computacionales y energéticos.

El Aprendizaje Federado (\textit{Federated Learning}, FL)~\cite{mcmahan2017federated} surgió como una alternativa descentralizada en la que múltiples nodos o clientes entrenan copias locales del modelo con sus propios datos privados, transmitiendo a un servidor central únicamente los parámetros o gradientes actualizados. Aunque FL solventa en gran medida la fuga de datos crudos, introduce una vulnerabilidad estructural: la falta de confianza intrínseca entre los participantes. En un entorno distribuido no supervisado, actores maliciosos o nodos comprometidos pueden ejecutar ataques de envenenamiento (\textit{data poisoning} o \textit{model poisoning}) inyectando pesos manipulados. Por ejemplo, en el contexto de los Vehículos Conectados y Autónomos (CAV), una pequeña alteración coordinada en los pesos podría inducir al modelo global a no clasificar correctamente una señal de tráfico crítica (como un STOP), con consecuencias potencialmente fatales.

Para mitigar este riesgo, la tecnología \textit{Blockchain} ofrece una solución idónea como registro contable inmutable, descentralizado y verificable~\cite{nakamoto2008bitcoin, zheng2018blockchain}. Al integrar Blockchain en un flujo de entrenamiento distribuido, cada actualización de parámetros, ronda de agregación e identidad de los clientes queda registrada criptográficamente mediante contratos inteligentes (\textit{smart contracts}), habilitando la trazabilidad absoluta y la no repudiación de las contribuciones.

Sin embargo, esta convergencia plantea un compromiso técnico fundamental: el dilema entre \textbf{seguridad y rendimiento}. Mientras que la Blockchain garantiza inmutabilidad y consenso bizantino, introduce una sobrecarga inevitable de latencia de red, tiempo de confirmación de transacciones y consumo de recursos. En aplicaciones de misión crítica o en tiempo cuasi-real, resulta imperativo responder a la siguiente pregunta: \textit{¿cuánta latencia y coste computacional estamos dispuestos a asumir a cambio de un entrenamiento auditable y más resistente a manipulaciones?}

\subsection{Motivación y Trabajos Previos}
Este trabajo se fundamenta y amplía dos líneas de investigación desarrolladas en la Universidad de Alicante:
\begin{enumerate}
    \item \textbf{Verificación de eventos vehiculares en tiempo real}~\cite{pujol2024blockchain}: demostró la viabilidad de una red privada Hyperledger Besu con consenso IBFT 2.0 para registrar alertas de tráfico con latencias de 0,32 a 0,82 segundos, mejorando el tiempo de respuesta humano.
    \item \textbf{Framework FL-Blockchain}~\cite{cloudlabaia_flblockchain}: planteó conceptualmente el uso de contratos inteligentes para registrar los \textit{hashes} criptográficos de los pesos de entrenamiento en cada ronda federada.
\end{enumerate}

La presente investigación valida, formaliza e implementa experimentalmente la convergencia de ambas propuestas en un banco de pruebas real, evaluando tanto la integridad como el impacto en el rendimiento temporal.

\subsection{Contribuciones Principales}
Las contribuciones de este artículo se resumen en los siguientes puntos:
\begin{itemize}
    \item \textbf{Diseño de una arquitectura híbrida}: Se formula un marco integral que acopla una red permisionada basada en Hyperledger Besu con el framework de aprendizaje federado Flower~\cite{beutel2020flower} y un sistema de almacenamiento distribuido direccionable por contenido (IPFS).
    \item \textbf{Trazabilidad e integridad mediante Smart Contracts}: Se diseña un contrato inteligente en Solidity que orquesta la verificación cruzada de modelos mediante resúmenes SHA-256, registrando de forma inalterable las contribuciones de cada cliente y el modelo global consolidado por ronda.
    \item \textbf{Evaluación empírica y mitigación de ataques adversariales}: Se somete la infraestructura a vectores de ataque deliberados de inyección de pesos corruptos (tanto en servidor comprometido como en clientes envenenadores), validando la capacidad del protocolo de doble verificación para detectar y neutralizar el fraude con una efectividad del 100\% y tiempos de reacción inferiores a $0{,}1$~segundos.
    \item \textbf{Evaluación cuantitativa de la infraestructura}: Se despliega un clasificador convolucional de tráfico con el propósito de emplear la estructura de aprendizaje federado como banco de pruebas experimental para extraer métricas cuantitativas reproducibles (sobrecoste temporal, consumo de gas y escalabilidad de validadores), delimitando la optimización fina del modelo como trabajo futuro.
    \item \textbf{Análisis de viabilidad operativa}: Se responde analíticamente al dilema latencia-seguridad, delimitando los escenarios en los que la sobrecarga del registro en cadena resulta plenamente asumible frente a esquemas no auditables.
\end{itemize}

\subsection{Estructura del Artículo}
El resto del manuscrito se estructura de la siguiente manera: la Sección~\ref{sec:estado_arte} expone el estado del arte y fundamentos teóricos de Blockchain y Aprendizaje Federado. La Sección~\ref{sec:arquitectura} detalla la arquitectura propuesta y el diseño del contrato inteligente. La Sección~\ref{sec:implementacion} describe el banco de pruebas experimental y el dataset empleado. La Sección~\ref{sec:evaluacion} presenta y discute los resultados de rendimiento y seguridad. Finalmente, la Sección~\ref{sec:conclusiones} resume las conclusiones y líneas futuras.

